\section{eTMF}

% SLIDE ID: sec07-goal
\BlockGoalFrame{\SlideID{sec07-goal}}{
  \item zrozumieć eTMF jako system gotowości dokumentacyjnej badania,
  \item zobaczyć, jak ważne są kompletność, terminowość i jakość dokumentów,
  \item połączyć eTMF z audit readiness i inspection readiness.
}

% SLIDE ID: sec07-tmf
\begin{frame}{Czym jest TMF}
  \SlideID{sec07-tmf}
  \begin{itemize}
    \item TMF to Trial Master File, czyli centralny zbiór dokumentów badania,
    \item dokumentuje planowanie, prowadzenie i nadzór nad badaniem,
    \item ma potwierdzać, że badanie było prowadzone zgodnie z wymaganiami.
  \end{itemize}
\end{frame}

% SLIDE ID: sec07-etmf
\begin{frame}{Czym jest eTMF}
  \SlideID{sec07-etmf}
  \begin{itemize}
    \item eTMF to elektroniczna wersja Trial Master File,
    \item porządkuje dokumenty, metadane, wersje i dostęp użytkowników,
    \item wspiera bieżącą gotowość do audytu i inspekcji.
  \end{itemize}
\end{frame}

% SLIDE ID: sec07-essential-docs
\begin{frame}{Essential documents}
  \SlideID{sec07-essential-docs}
  \begin{itemize}
    \item protokół i jego wersje,
    \item zgody, umowy, listy delegacji, potwierdzenia szkoleń,
    \item dokumenty regulatorowe, safety i monitoringowe,
    \item korespondencja i artefakty wymagane przez model badania.
  \end{itemize}
\end{frame}

% SLIDE ID: sec07-structure
\begin{frame}{Struktura dokumentacji}
  \SlideID{sec07-structure}
  \begin{itemize}
    \item dokumenty są porządkowane według obszaru, typu i etapu badania,
    \item ważne są nie tylko pliki, ale też metadane i status dokumentu,
    \item dobra struktura skraca czas odnalezienia właściwego materiału.
  \end{itemize}
\end{frame}

% SLIDE ID: sec07-tmf-reference-model-what
\begin{frame}{Trial Master File Reference Model}
  \SlideID{sec07-tmf-reference-model-what}
  \begin{itemize}
    \item to branżowy model referencyjny porządkujący zawartość \texttt{TMF} i \texttt{eTMF},
    \item projekt rozpoczął się w 2008 roku, a pierwsza wersja została opublikowana w 2010 roku,
    \item jego celem jest ujednolicenie nazewnictwa, struktury i metadanych między sponsorami, CRO i ośrodkami,
    \item model nie jest narzuconą jedyną strukturą, ale punktem odniesienia dla harmonizacji i inspection readiness.
  \end{itemize}
\end{frame}

% SLIDE ID: sec07-tmf-reference-model-version
\begin{frame}{Aktualna wersja i znaczenie praktyczne}
  \SlideID{sec07-tmf-reference-model-version}
  \begin{itemize}
    \item na stronie \texttt{CDISC} aktualnie wskazana jest wersja \texttt{3.3.1}, opublikowana 15 sierpnia 2023 roku,
    \item wcześniejsze wersje rozwijały model od \texttt{v1.0} przez \texttt{v2.0}, \texttt{v3.0}, \texttt{v3.1}, \texttt{v3.2.0} i \texttt{v3.3},
    \item aktualizacje odzwierciedlają zmiany regulacyjne, praktykę inspekcyjną i potrzeby pracy w systemach elektronicznych,
    \item w praktyce \texttt{eTMF} często mapuje własną konfigurację właśnie do tego modelu.
  \end{itemize}
\end{frame}

% SLIDE ID: sec07-tmf-reference-model-structure
\begin{frame}{Z czego składa się TMF Reference Model}
  \SlideID{sec07-tmf-reference-model-structure}
  \begin{itemize}
    \item model ma strukturę: \textbf{Zone → Section → Artifact}, a w nowszych wersjach także \textbf{sub-artifacts},
    \item oficjalnie obejmuje \textbf{11 stref}, m.in. Trial Management, Central Trial Documents, Regulatory, Site Management, Safety Reporting, Data Management i Statistics,
    \item każdy artifact ma nazwę, definicję/purpose, poziom \texttt{core} lub \texttt{recommended}, odniesienia regulacyjne i identyfikator,
    \item to ważne, bo \texttt{eTMF} nie zarządza tylko plikami, ale też logiką klasyfikacji i wyszukiwania dokumentów.
  \end{itemize}
\end{frame}

% SLIDE ID: sec07-roles
\begin{frame}{Role i odpowiedzialności}
  \SlideID{sec07-roles}
  \begin{itemize}
    \item ktoś odpowiada za wytworzenie dokumentu,
    \item ktoś odpowiada za upload i klasyfikację,
    \item ktoś przegląda kompletność i jakość,
    \item odpowiedzialność musi być jasna po stronie sponsora i zespołu badania.
  \end{itemize}
\end{frame}

% SLIDE ID: sec07-ctq
\begin{frame}{Completeness, timeliness, quality}
  \SlideID{sec07-ctq}
  \begin{itemize}
    \item \textbf{completeness}: czy wszystko, co powinno być, rzeczywiście jest,
    \item \textbf{timeliness}: czy dokument trafił do systemu we właściwym czasie,
    \item \textbf{quality}: czy dokument jest poprawny, aktualny i właściwie opisany.
  \end{itemize}
\end{frame}

% SLIDE ID: sec07-audit-inspection
\begin{frame}{Audit readiness i inspection readiness}
  \SlideID{sec07-audit-inspection}
  \begin{itemize}
    \item audit readiness oznacza gotowość do wewnętrznego lub zewnętrznego przeglądu,
    \item inspection readiness oznacza gotowość do kontroli regulatora,
    \item eTMF ma umożliwiać szybkie pokazanie kompletnych i wiarygodnych dokumentów.
  \end{itemize}
\end{frame}

% SLIDE ID: sec07-versioning
\begin{frame}{Wersjonowanie, podpisy i dostęp}
  \SlideID{sec07-versioning}
  \begin{itemize}
    \item wersjonowanie pozwala odróżnić dokument aktualny od archiwalnego,
    \item podpisy elektroniczne wspierają zatwierdzanie i ścieżkę odpowiedzialności,
    \item kontrola dostępu ogranicza widoczność i edycję do uprawnionych ról.
  \end{itemize}
\end{frame}

% SLIDE ID: sec07-problems
\begin{frame}{Najczęstsze problemy w eTMF}
  \SlideID{sec07-problems}
  \begin{itemize}
    \item dokument jest w systemie, ale pod złą kategorią,
    \item dokument jest nieaktualny albo bez właściwej wersji,
    \item dokument pojawia się z opóźnieniem,
    \item nie wiadomo, kto odpowiada za brak lub poprawkę.
  \end{itemize}
\end{frame}

% SLIDE ID: sec07-case
\CaseStudyFrame{Inspekcja zaczyna się jutro — czego szukamy w eTMF?}{
  Trzeba szybko ocenić, czy dokumentacja badania jest kompletna, aktualna i gotowa do okazania audytorowi lub inspektorowi.
}{
  Szukamy braków, opóźnionych uploadów, błędnych wersji, niewłaściwych metadanych, nieczytelnych podpisów oraz dokumentów bez jasnego właściciela.
}

% SLIDE ID: sec07-discussion
\DiscussionFrame{Dlaczego eTMF nie jest tylko folderem z plikami?}{
  Które elementy sprawiają, że \texttt{eTMF} jest systemem jakości i gotowości inspekcyjnej, a nie zwykłym repozytorium? Co się psuje, jeśli dokument jest wgrany, ale bez metadanych, wersji albo właściciela?
}

% SLIDE ID: sec07-sources
\begin{frame}{Źródła do bloku eTMF}
  \SlideID{sec07-sources}
  \begin{itemize}
    \item \texttt{CDISC Trial Master File Reference Model}: opis standardu i lista wersji, w tym \texttt{v3.3.1} z 15.08.2023,
    \item \texttt{TMF Reference Model Forum}: historia projektu, pierwsza publikacja w 2010 roku i rozwój od inicjatywy rozpoczętej w 2008 roku,
    \item \texttt{TMF Reference Model User Guide}: struktura \texttt{Zone → Section → Artifact} oraz 11 stref modelu,
    \item \texttt{ICH E6 GCP} i wytyczne \texttt{TMF}: kontekst regulacyjny dla essential documents i gotowości inspekcyjnej.
  \end{itemize}
\end{frame}
